(a) Let $\zeta$ be the generic point of $Z$, and choose a system of parameters of $\mathcal O_{Y,\eta'}$. It has $\operatorname{codim}(Y',Y)$ elements, and its image in $\mathcal O_{X,\zeta}$ generates an ideal whose radical is the maximal ideal, since $Z$ is a component of $f^{-1}(Y')$. The dimension theorem for noetherian local rings gives $\dim\mathcal O_{X,\zeta}\le\dim\mathcal O_{Y,\eta'}$ (Atiyah--Macdonald, \emph{Introduction to Commutative Algebra}, Theorem 11.14). Now use \exref{II.3.20}(c).

(b) Let $W$ be an irreducible component of $X_y$, and give its closure $Z$ in $X$ the reduced structure. Then $Z$ is a component of $f^{-1}(\overline{\{y\}})$ dominating $\overline{\{y\}}$, and $W$ is its generic fiber. By \exref{II.3.20}(b), additivity of transcendence degrees, and (a),
\[\dim W=\dim Z-\dim\overline{\{y\}}\ge\dim X-\dim Y=e.\]

(c) Restrict to nonempty affine open subsets $X=\operatorname{Spec}A$ and $Y=\operatorname{Spec}B$. Choose $t_1,\ldots,t_e\in A$ algebraically independent over $K(Y)$ such that $K(X)$ is finite over $K(Y)(t_1,\ldots,t_e)$. The map $X\to\operatorname{Spec}B[t_1,\ldots,t_e]=\mathbb A_Y^e$ is generically finite. By \exref{II.3.7}, it is finite over a nonempty open subset $W$ of $\mathbb A_Y^e$. The resulting finite map $U\to W$ is surjective, since it is dominant. Its nonempty fibers over $Y$ are finite surjective over the corresponding open subsets of $\mathbb A_{k(y)}^e$, so each has dimension $e$.

(d) Part (b) gives $E_h=X$ for $h\le e$. For $h>e$, choose $U$ as in (c). A component of a fiber of dimension at least $h$ cannot meet $U$, since a nonempty open subset of an integral scheme of finite type over a field has the same dimension. Thus $E_h\subseteq X\setminus U$, proving (2). Let $Z_1,\ldots,Z_r$ be the reduced irreducible components of $X\setminus U$. For each $i$, apply the same definition of $E_h$ to the dominant morphism $Z_i\to\overline{f(Z_i)}$. The union of these sets equals $E_h$: every large component of a fiber of $X$ lies in some $Z_i$, and every large component of a fiber of $Z_i$ lies in a component of the corresponding fiber of $X$ of at least the same dimension. Each set is closed by induction on $\dim X$, hence $E_h$ is closed.

(e) For $h\ge0$, $C_h=f(E_h)\setminus f(E_{h+1})$, so it is constructible by \exref{II.3.19}. The generic fiber has dimension $e$ by \exref{II.3.20}(b), so the generic point of $Y$ belongs to $C_e$. A constructible subset of an irreducible noetherian space containing its generic point contains a dense open subset by \exref{II.3.18}. The empty-fiber locus is $Y\setminus f(X)$, also constructible.
